% !TeX program = xelatex
% 2026-10-06针对性文字精修；保留原稿实验数值、公式、表格及引用键。
\documentclass[UTF8,a4paper,12pt,fontset=fandol]{ctexart}
\usepackage[margin=2.5cm]{geometry}
\usepackage{amsmath,amssymb}
\usepackage{booktabs,longtable,tabularx,array,calc}
\usepackage{graphicx,xcolor}
\usepackage{placeins}
\usepackage{etoolbox,needspace}
\usepackage{hyperref}
\hypersetup{colorlinks=true,linkcolor=black,citecolor=black,urlcolor=blue}
\setlength{\parindent}{2em}
\setlength{\parskip}{0.25em}
\setlength{\emergencystretch}{3em}
\linespread{1.22}
\renewcommand{\arraystretch}{1.22}
\AtBeginEnvironment{longtable}{\small\setlength{\tabcolsep}{3pt}}
\providecommand{\tightlist}{\setlength{\itemsep}{0pt}\setlength{\parskip}{0pt}}
\setcounter{secnumdepth}{0}
\title{预训练原子势微调中元素参考能差异对相对能量的影响\\[0.45em]\large Energy-reference alignment reduces excess structure-dependent changes during energy-only adaptation of pretrained interatomic potentials}
\author{}
\date{}
\begin{document}
\maketitle

\begin{abstract}
预训练机器学习原子势可以通过有限的新数据进一步适配，但不同数据之间的元素参考能差异可能影响这一过程。在只有能量标签的情况下，如果这部分差异直接由模型通过微调拟合，绝对能量误差虽然能够降低，结构间原有的相对能量和原子力也可能同时发生额外变化。
本文以 SevenNet 为主要模型，比较直接微调与训练前先处理元素参考能差异两种方式，并通过化学组成相同的结构之间的相对能量判断模型实际产生的变化是否符合目标要求。结果表明，直接微调在降低形成能误差的同时产生了较多目标要求之外的相对能量变化，而提前处理元素参考能后这些变化明显减少。构造的受控实验只改变元素参考能，而保持目标相对能量和原子力不变，直接微调仍产生了明显的额外变化，提前处理元素参考能后则基本回到正常训练的背景水平。在 EquiformerV3 和 GGA/GGA+\(U\) 到 r$^2$SCAN 的适配中也观察到相同趋势。
这些结果表明，在仅用能量标签微调中，元素参考能差异不应简单交由模型参数更新吸收。提前处理这部分差异能够减少目标并不需要的结构间相对能量和原子力变化，使微调更集中于目标数据真正要求的能量变化。
\end{abstract}
\noindent\textbf{关键词：}机器学习原子势；仅能量适配；能量参考对齐；形成能；相对能量

\section{引言}
材料的热力学稳定性是新材料发现中的重要问题，也是判断候选结构能否实际存在和合成的重要依据\cite{sun2016thermodynamic}，而形成能反映了化合物相对于元素参考态的能量变化，是构建相稳定性凸包和判断材料热力学稳定性的基础\cite{bartel2020critical}。传统形成能通常依赖密度泛函理论（DFT）计算，在大规模材料筛选中计算成本较高\cite{curtarolo2013high}。随着材料数据不断积累，机器学习逐渐被用于形成能预测\cite{schmidt2019recent}。早期的 CGCNN、MEGNet 和 ALIGNN 等模型主要以晶体结构为输入，预测形成能、带隙等材料性质\cite{xie2018crystal,chen2019graph,choudhary2021atomistic}；随后，SchNet、DimeNet 等模型开始直接利用原子坐标学习体系能量及其对结构的依赖，使机器学习模型能够进一步用于势能面建模和原子尺度模拟\cite{schutt2017schnet,gasteiger2020directional}。NequIP、MACE 等等变模型进一步提高了对原子局部环境和相互作用的描述能力\cite{batzner20223,batatia2022mace}。随着训练数据和材料体系不断扩大，M3GNet、CHGNet、MACE-MP、SevenNet 和 EquiformerV3 等模型开始在大量不同元素和结构上进行训练，使一个模型能够覆盖更广的材料范围，并作为后续新体系和新能量数据适配的起点\cite{chen2022universal,deng2023chgnet,batatia2025foundation,park2024scalable,liao2026equiformerv3}。

与此同时，大规模第一性原理数据库和数据平台也为这些模型的训练提供了数据基础\cite{batra2021emerging}。Materials Project、OQMD、AFLOW、JARVIS-DFT、NOMAD 和 Alexandria 等数据库积累了大量晶体结构及其能量、形成能、电子结构和其他材料性质数据，覆盖了较广的元素组合和材料体系。其中，Materials Project 和 OQMD 被广泛用于形成能和热力学稳定性研究\cite{jain2013commentary,kirklin2015open},，JARVIS-DFT 和 AFLOW 提供了较丰富的结构与多种材料性质数据\cite{curtarolo2012aflowlib,choudhary2020joint}，NOMAD 和 Alexandria 则进一步扩展了可用于机器学习训练的大规模第一性原理计算数据\cite{draxl2018nomad,schmidt2024improving}。

然而，预训练模型通常基于一个或多个数据库进行训练。不同数据库采用的计算方法、参数设置和元素参考能并不完全相同，因此同一种材料在不同数据中的能量可能存在差异\cite{hegde2023quantifying}。当模型进一步用于新的数据库或新的计算数据时，这些差异会直接反映在形成能预测误差中，通常需要通过微调进行适配\cite{huang2025cross}。问题在于，模型只能看到预测值与目标值之间的误差，却不能区分其中哪些来自元素参考能不同，哪些来自原子排列所对应的真实能量差异。如果这些误差全部由模型在微调中吸收，形成能误差虽然可以降低，但原有的不同构型之间的能量关系也可能同时发生改变。

这一问题在仅有能量标签的适配中尤其值得关注。在部分高精度电子结构计算中，高精度能量相对容易获得，而相应的原子力或能量梯度计算更为困难\cite{messerly2025multi}，因此已有工作采用有限的高精度能量对预训练模型进行适配\cite{zaverkin2023transfer}。当目标数据只有能量标签时，训练只能直接约束各个结构的预测能量，不能直接约束能量随原子位置变化的方式\cite{christensen2020role}。即使形成能误差明显降低，也不能保证不同构型之间的相对能量和原子力得到正确保持。仅比较微调前后的绝对能量 MAE，也无法判断模型实际改变了哪些能量关系，以及这些变化中哪些是目标数据真正需要的，哪些可能与元素参考能差异有关。

要区分这些变化，可以利用元素参考能的一个基本特点。元素参考能只与材料中元素的种类和数量有关，而与原子的具体排列无关\cite{friedrich2019coordination}。因此，对于化学组成相同、但原子排列不同的结构，元素参考能带来的共同偏移会在比较能量差时抵消。这些不同构型之间的相对能量可以用来判断微调是否改变了原有的能量关系。同时，元素参考能不依赖原子坐标，本身也不会改变原子力。如果目标本身不要求改变这些构型之间的相对能量和原子力，而微调后它们仍然发生了变化，就说明模型还产生了目标之外的改变。

元素参考能带来的能量差异，在材料计算中早已受到关注。在形成能计算中，人们常通过重新确定元素参考能来减小不同计算方法或实验数据之间的系统偏差\cite{stevanovic2012correcting}。随着预训练原子势开始用于不同计算方法和数据库之间的迁移，这类差异也进入了模型适配过程。Huang 等人的工作表明，在从 GGA/GGA+$U$ 向 r$^2$SCAN 迁移时，重新确定目标数据的元素参考能可以明显改善模型训练和预测结果\cite{huang2025cross}；Kim 等人进一步指出，在多个数据库之间进行训练和迁移时，不同数据的原子能量参考也需要分别处理\cite{kim2026optimizing}。这些工作已经说明，正确处理元素参考能对模型迁移十分重要。但如果这部分差异没有提前处理，而是直接让网络在仅有能量标签的微调中拟合，模型是否会同时改变本来不需要变化的构型间相对能量和原子力，以及这些改变与目标数据真正要求的变化有多大差别，仍缺少直接分析。

本文从化学组成相同、原子排列不同的结构入手，考察微调是否改变了它们之间的相对能量。通过比较目标数据实际要求的能量变化与微调后模型产生的变化，判断模型是否进行了目标之外的调整。我们首先以 SevenNet 对 Matbench 形成能数据进行适配，比较直接微调和先处理元素参考能差异后再微调两种方式。直接微调将形成能 MAE 降至 (69.9~\mathrm{meV/atom})，但同时明显改变了不同构型之间的相对能量；提前处理这部分差异后，MAE 进一步降至 (21.5~\mathrm{meV/atom})，模型产生的额外变化也明显减小。
为了单独检验元素参考能差异本身的影响，我们又构造了一个只需要修正绝对能量、而相对能量和原子力都不应发生变化的受控任务。即使在这种情况下，直接微调仍然产生了明显的相对能量和原子力变化，而提前处理元素参考能后，这些变化接近正常训练的背景水平。相同趋势也在 EquiformerV3 以及 GGA/GGA+\(U\) 到 r$^2$SCAN的跨泛函适配中得到验证。

这些结果表明，在仅用能量标签的微调中，形成能误差的降低并不意味着模型只学习了目标真正需要的变化。通过比较目标要求的变化与模型实际产生的变化，可以识别微调过程中出现的额外调整，而提前处理元素参考能差异能够明显减少这类不必要的变化。

\section{结果}

\subsection{提前处理元素参考能差异可减少微调对结构间相对能量的改变}

我们选择预训练机器学习原子势 SevenNet 作为研究对象，并使用 Matbench 中的 Materials Project 形成能数据作为目标数据。该数据集提供了统一整理的晶体结构和形成能标签，适合用于比较不同微调方法。我们比较两种训练方式：第一种直接在这套形成能数据上继续训练，记为 direct；第二种先利用训练集拟合新数据的元素参考能，再进行训练，记为 aligned。为了判断拟合参考能本身能够带来多少变化，我们还设置了两个不训练的对照：一个是原始预训练模型 Frozen，另一个是在 Frozen 的预测结果上只加入前面拟合出的参考能修正，记为 Frozen+$h$。direct 和 aligned 使用相同的数据划分和训练设置。训练结束后，如果预测中仍然存在只与元素组成有关的系统误差，两种方法都采用相同的方式进行修正，以保证比较公平。

我们从两个方面评价模型的适配结果。首先使用形成能 MAE 衡量模型对绝对形成能的预测误差；随后，对于化学组成相同但原子排列不同的结构，比较它们之间的能量差，并统计模型能否正确判断每组中能量最低的结构。前者反映模型对形成能整体大小的预测能力，后两项则反映微调后不同结构之间的相对能量关系是否保持正确。表1给出三个固定随机种子的结果。

\textbf{表1｜SevenNet 在不同训练方式下的形成能与结构间相对能量结果。}

形成能 MAE 和结构间能量差 MAE 的单位均为 meV/atom；最低能结构命中率表示模型正确找出每组最低能结构的比例。direct 和 aligned 的``$\pm$''表示三个固定随机种子的标准差。

\begin{longtable}[]{@{}
  >{\raggedright\arraybackslash}p{0.22\columnwidth}
  >{\raggedleft\arraybackslash}p{0.25\columnwidth}
  >{\raggedleft\arraybackslash}p{0.28\columnwidth}
  >{\raggedleft\arraybackslash}p{0.20\columnwidth}@{}}
\toprule
模型 & 形成能 MAE & 结构间能量差 MAE & 最低能结构命中率 \\
\midrule
\endhead
\bottomrule
\endlastfoot

Frozen
& 4484.523
& 22.014
& 0.863 \\

Frozen+$h$
& 109.675
& 22.014
& 0.863 \\

direct
& 69.911 $\pm$ 0.503
& 67.171 $\pm$ 0.563
& 0.719 $\pm$ 0.008 \\

aligned
& 21.474 $\pm$ 0.160
& 23.379 $\pm$ 0.100
& 0.844 $\pm$ 0.006 \\

\end{longtable}

原始 SevenNet 的形成能 MAE 高达 \(4484.523~\mathrm{meV/atom}\)，但对于化学组成相同的不同结构，其能量差 MAE 只有 \(22.014~\mathrm{meV/atom}\)，最低能结构命中率达到 0.863。只加入前面拟合的元素参考能修正后，形成能 MAE 直接降至 \(109.675~\mathrm{meV/atom}\)，而结构间能量差和最低能结构命中率完全不变。这说明预训练模型与目标形成能之间很大一部分误差可以通过元素参考能的修正消除，并不需要改变不同结构之间原有的能量关系。

训练结束后，两种微调方式出现了明显差别。direct 将形成能 MAE 进一步降至 \(69.911~\mathrm{meV/atom}\)，但结构间能量差 MAE 增加到 \(67.171~\mathrm{meV/atom}\)，最低能结构命中率也由 0.863 降至 0.719。相比之下，aligned 的形成能 MAE 降至 \(21.474~\mathrm{meV/atom}\)，结构间能量差 MAE 为 \(23.379~\mathrm{meV/atom}\)，最低能结构命中率为 0.844，均接近微调前的水平。也就是说，直接微调虽然降低了形成能误差，却同时明显改变了不同结构之间的相对能量；提前处理元素参考能差异后，形成能预测更加准确，同时减少了这种额外改变。

表1说明 direct 和 aligned 对不同结构之间相对能量的影响明显不同，但微调以后模型和 Frozen 不一样，并不能马上说明模型改错了，因为目标数据里的真实相对能量本来就可能和 Frozen 不一样。我们进一步比较目标数据实际要求的变化与微调后模型真正产生的变化。

对于化学组成相同、原子排列不同的一对结构 \(a\) 和 \(b\)，预训练模型给出的单位原子能量记为 \(e_0\)，目标能量记为 \(y\)，微调后的预测记为 \(e_s\)。目标数据要求这两个结构之间的能量差发生

\[
t_{ab}
=
(y_a-y_b)
-
(e_{0,a}-e_{0,b})
\]

的变化，而模型经过微调后实际产生的变化为

\[
c_{s,ab}
=
(e_{s,a}-e_{s,b})
-
(e_{0,a}-e_{0,b}).
\]

如果 \(c_{s,ab}=t_{ab}\)，说明模型对这一对结构的调整与目标完全一致。由于两个结构具有相同的元素组成，元素参考能带来的共同偏移会在能量差中抵消，因此这里比较的是不同结构之间的相对能量变化。

我们首先比较模型产生的总体变化与目标要求的变化：

\[
\frac{\|c_s\|}{\|t\|}.
\]

比值大于1表示模型总体改变得更多，小于1表示改变得较少。但变化幅度相同并不代表每一对结构都改对了，因此我们进一步把模型产生的变化分成沿目标变化方向和偏离目标方向的两部分。沿目标方向的比例定义为

\[
\alpha_s
=
\frac{c_s^{\mathsf T}t}{t^{\mathsf T}t},
\]

剩余的

\[
\frac{\|c_s-\alpha_s t\|}{\|t\|}
\]

表示模型在目标要求之外产生的变化。理想情况下，模型实际产生的变化与目标完全一致，此时总体变化比和 \(\alpha_s\) 均为1，目标方向之外的变化为0。

\Needspace{10\baselineskip}
\textbf{表2｜SevenNet 微调产生的结构间能量差变化与目标要求的比较。}

\begin{longtable}[]{@{}
  >{\raggedright\arraybackslash}p{0.22\columnwidth}
  >{\raggedleft\arraybackslash}p{0.24\columnwidth}
  >{\raggedleft\arraybackslash}p{0.24\columnwidth}
  >{\raggedleft\arraybackslash}p{0.26\columnwidth}@{}}
\toprule
训练方式 &
总体变化／目标变化 &
沿目标方向的比例 &
目标方向之外的变化／目标变化 \\
\midrule
\endhead
\bottomrule
\endlastfoot

direct &
2.641 $\pm$ 0.016 &
0.498 $\pm$ 0.026 &
2.593 $\pm$ 0.021 \\

aligned &
0.741 $\pm$ 0.007 &
0.358 $\pm$ 0.000 &
0.649 $\pm$ 0.008 \\

\end{longtable}

direct 产生的总体变化达到目标要求的 2.641 倍，沿目标方向的比例为 0.498，而目标方向之外的变化达到目标要求的 2.593 倍。这说明 direct 并不是简单地把正确的变化放大了，而是在微调过程中产生了大量目标没有要求的调整。

aligned 产生的总体变化为目标要求的 0.741 倍，目标方向之外的变化降至 0.649。虽然 aligned 同样没有完全实现目标数据要求的结构间能量变化，但与 direct 相比，它产生的额外调整明显更少。

我们还直接比较了模型实际变化与目标变化之间的距离：

\[
\frac{\|c_s-t\|}{\|t\|}.
\]

如果模型完全不改变预训练结果，该值为1。direct 和 aligned 分别为 2.641 和 0.913，说明 direct 产生的结构间能量变化比保持原模型不变更偏离目标，而 aligned 明显减少了这种偏离。配对 bootstrap 结果也表明，aligned 与 direct 之间的差异明确。

为了确认两种训练方式的差别不是由具体训练设置造成的，我们改变了学习率，并测试在微调过程中继续调整元素参考能。不同设置下，direct 与 aligned 之间的主要差距仍然存在，说明前面的结果并不是某一学习率或参考能在训练中是否继续调整造成的。训练记录同时显示，direct 在训练过程中面对的残差更大，网络参数的整体变化也更多。

综合以上结果，直接微调虽然能够降低形成能误差，但同时产生了大量目标数据并不要求的结构间能量变化。提前处理元素参考能差异后，绝对形成能预测进一步改善，这些额外变化也明显减少。但在真实形成能数据中，目标本身可能确实要求部分结构间能量发生变化，因此还不能仅凭这一结果判断这些额外变化是不是由元素参考能差异造成的。

\subsection{只改变元素参考能也会引起额外的相对能量和原子力变化}

为了检查元素参考能差异本身会不会带来这些额外变化，我们设计了一个更简单的受控实验：只改变元素参考能，而不改变目标中的结构间相对能量和原子力。

我们以冻结 SevenNet 的预测能量 \(E_0(X)\) 为起点，在其基础上加入训练集拟合得到的元素参考能 \(h(z)\)：

\[
E_{\mathrm{syn}}(X)
=
E_0(X)+N_Xh(z),
\]

其中 \(N_X\) 为结构中的原子数。\(h(z)\) 只与元素组成有关，因此对于化学组成相同、原子排列不同的结构，这一项完全相同，它们之间的能量差不会改变。同时，\(h(z)\) 不依赖原子坐标，因此

\[
\nabla_{\mathbf R}[N_Xh(z)]=0,
\]

加入这一项也不会改变目标原子力，所以这个任务真正要求模型完成的只有绝对能量上的参考能变化。

我们设置三种训练方式。第一种直接重新拟合原始预训练能量，记为 zero\_direct，用来测量正常训练本身会带来多大的变化；第二种直接拟合加入元素参考能后的目标，记为 reference\_direct，用来观察当这部分差异完全交给网络学习时会发生什么；第三种使用相同的目标，但在训练前先加入这部分元素参考能，记为 reference\_aligned。三种训练使用相同的训练设置、训练轮数和随机种子。

\textbf{表3｜只改变元素参考能时的受控实验结果。}

结构间能量差变化表示训练前后相同化学组成结构之间能量差的变化，单位为 meV/atom；原子力变化为训练后相对于预训练模型的力分量 RMSE，单位为 eV/\AA。

\begin{longtable}[]{@{}
  >{\raggedright\arraybackslash}p{0.28\columnwidth}
  >{\raggedleft\arraybackslash}p{0.22\columnwidth}
  >{\raggedleft\arraybackslash}p{0.25\columnwidth}
  >{\raggedleft\arraybackslash}p{0.22\columnwidth}@{}}
\toprule
训练方式 &
绝对能量 MAE &
结构间能量差变化 &
原子力变化 RMSE \\
\midrule
\endhead
\bottomrule
\endlastfoot

zero\_direct &
0.902 $\pm$ 0.052 &
0.207 $\pm$ 0.023 &
0.00415 $\pm$ 0.00037 \\

reference\_direct &
53.990 $\pm$ 0.046 &
18.770 $\pm$ 0.304 &
0.74877 $\pm$ 0.00957 \\

reference\_aligned &
0.995 $\pm$ 0.052 &
0.210 $\pm$ 0.041 &
0.00382 $\pm$ 0.00087 \\

\end{longtable}

结果很明显。zero\_direct 中，结构间能量差和原子力都只发生了很小的变化，说明正常训练本身会带来少量背景漂移。加入元素参考能差异后，如果直接让网络拟合，结构间能量差变化增大到 \(18.770~\mathrm{meV/atom}\)，原子力变化达到 \(0.74877~\mathrm{eV/\AA}\)，远高于 zero\_direct 的背景水平。而在这个受控任务中，新增的目标只要求修正绝对能量，并不要求结构间相对能量和原子力发生变化。提前处理元素参考能后，对应变化分别降至 \(0.210~\mathrm{meV/atom}\) 和 \(0.00382~\mathrm{eV/\AA}\)，重新回到接近 zero\_direct 的水平。说明当元素参考能差异直接交给网络拟合时，微调会同时带来目标没有要求的相对能量和原子力变化。

这一现象在扩大原子力评价范围后仍然存在。reference\_direct 的平均原子力变化为 \(0.4169~\mathrm{eV/\AA}\)，而 reference\_aligned 和 zero\_direct 分别只有 \(0.00247\) 和 \(0.00235~\mathrm{eV/\AA}\)。在评价的 100 个结构中，direct 的原子力变化均更大，说明这一结果并不是由少量特殊结构造成的。

还需要排除另一种可能：reference\_direct 出现较大变化，也许只是因为它没有把目标中的元素参考能学好。为检查这一点，我们比较模型实际产生的能量变化 \(u\) 与人工加入的元素参考能变化 \(h\)，计算

\[
\beta
=
\frac{u^{\mathsf T}h}{h^{\mathsf T}h}.
\]

如果模型已经基本学到这部分元素参考能变化，$\beta$ 应接近 1。对于 reference\_direct，$\beta=0.998613\pm0.001182$；对于 reference\_aligned，$\beta=1.000052\pm0.000008$。两者均接近 1，说明 direct 实际上已经学到了目标中的元素参考能变化。

进一步扣除这部分参考能变化后，我们用 $\|u-\beta h\|/\|h\|$ 衡量剩余能量变化的大小。reference\_direct 的结果为 $0.022250\pm0.000140$，而 reference\_aligned 仅为 $0.000481\pm0.000045$。因此，direct 的主要问题并不是没有学到元素参考能，而是在完成这部分能量变化的同时，还引入了更多与参考能变化无关的额外变化。

综合以上结果，在一个只要求修正元素参考能的任务中，直接微调仍然会明显改变结构间相对能量和原子力；提前处理元素参考能后，这些变化基本回到正常训练的背景水平。这说明，即使目标只涉及元素参考能变化，直接进行仅能量标签微调时，模型仍可能同时改变本来不需要变化的结构间能量关系和原子力。

\subsection{提前处理元素参考能差异的效果可在其他模型和能量数据中复现}

前面的结果主要来自 SevenNet。为了判断这一现象是否依赖具体模型，我们进一步在 EquiformerV3（EqV3）上比较 direct 和 aligned,还测试了从 GGA/GGA+\(U\) 能量向 r$^2$SCAN 能量进行适配的任务，来判断更换目标能量计算方法后是否仍有相同趋势。每项任务中，direct 和 aligned 均使用相同的初始模型、数据划分和训练设置。
\textbf{表4｜EqV3 和 GGA/GGA+$U$$\rightarrow$r$^2$SCAN 任务中的 direct/aligned 比较。}

形成能 MAE 和结构间能量差 MAE 的单位均为 meV/atom；最低能结构命中率表示模型正确找出每组最低能结构的比例。不同任务的数据集和评价集合不同，因此表中的数值只用于比较同一任务内的 direct 和 aligned。

\begin{longtable}[]{@{}
  >{\raggedright\arraybackslash}p{0.22\columnwidth}
  >{\raggedright\arraybackslash}p{0.18\columnwidth}
  >{\raggedleft\arraybackslash}p{0.20\columnwidth}
  >{\raggedleft\arraybackslash}p{0.22\columnwidth}
  >{\raggedleft\arraybackslash}p{0.18\columnwidth}@{}}
\toprule
模型／任务 &
训练方式 &
形成能 MAE &
结构间能量差 MAE &
最低能结构命中率 \\
\midrule
\endhead
\bottomrule
\endlastfoot

EqV3 &
direct &
62.384 $\pm$ 1.681 &
54.569 $\pm$ 1.314 &
0.684 $\pm$ 0.027 \\

EqV3 &
aligned &
17.789 $\pm$ 0.455 &
17.154 $\pm$ 1.958 &
0.840 $\pm$ 0.013 \\

GGA/GGA+$U$$\rightarrow$r$^2$SCAN &
direct &
83.667 $\pm$ 0.114 &
24.175 $\pm$ 0.662 &
0.463 $\pm$ 0.006 \\

GGA/GGA+$U$$\rightarrow$r$^2$SCAN &
aligned &
23.943 $\pm$ 0.092 &
12.979 $\pm$ 0.071 &
0.744 $\pm$ 0.000 \\

\end{longtable}

两项任务都得到与 SevenNet 相同的趋势。在 EqV3 中，提前处理元素参考能差异后，形成能 MAE 由 \(62.384\) 降至 \(17.789~\mathrm{meV/atom}\)，结构间能量差 MAE 由 \(54.569\) 降至 \(17.154~\mathrm{meV/atom}\)，最低能结构命中率也由 0.684 提高到 0.840。在 GGA/GGA+$U$$\rightarrow$r$^2$SCAN 的任务中，形成能 MAE 由 \(83.667\) 降至 \(23.943~\mathrm{meV/atom}\)，结构间能量差 MAE 由 \(24.175\) 降至 \(12.979~\mathrm{meV/atom}\)，最低能结构命中率由 0.463 提高到 0.744。说明前面观察到的差别并不只出现在 SevenNet 中，更换模型和目标能量数据后仍然存在。

我们还进一步检查了微调前后已有结构之间的能量关系。为避免重新弛豫结构带来其他影响，我们预先固定了一组 Materials Project 的弛豫轨迹，共包含 50 个计算任务、171 个去重结构和 434 对同一任务中的结构，比较 Frozen、direct 和 aligned 对这些固定结构给出的相对能量。

逐任务计算结构间能量差误差后，其中位数在 Frozen、aligned 和 direct 下分别为 \(54.4\)、\(84.9\) 和 \(137.3~\mathrm{meV/atom}\)。在 50 个任务中，aligned 有 41 个任务优于 direct，说明提前处理元素参考能差异能够明显减轻直接微调造成的相对能量变化。对异常结构以及不同统计方式进行检查后，这一趋势仍然保持。

综合 SevenNet、EqV3、GGA/GGA+$U$$\rightarrow$r$^2$SCAN 以及固定弛豫轨迹的结果，提前处理元素参考能差异带来的改善并不依赖某一个模型或某一套目标能量数据。其主要作用是减少直接微调过程中额外产生的结构间相对能量变化。

\section{讨论}

本文结果表明，仅用能量标签微调虽然可以降低目标形成能误差，但也可能同时改变预训练模型已经学到的不同结构之间的相对能量。在本文任务中，我们将预训练 SevenNet 进一步适配到 Matbench 的 Materials Project 形成能标签。由于预训练模型原有的能量输出与这套形成能标签在定义和元素参考能上存在差异，SevenNet 未经适配时具有很大的绝对形成能误差。只加入由训练集拟合得到的元素参考能后，绝对形成能误差已经大幅下降，而不同结构之间的能量差没有改变。这说明预训练模型与目标数据之间的一部分误差并不需要通过改变模型原有的结构信息来解决。进一步微调时，direct 虽然继续降低了形成能误差，却产生了较大的额外相对能量变化；提前处理元素参考能差异后，这些变化明显减少。

这种差别与两种训练方式面对的误差不同有关。direct 直接拟合预训练预测与目标标签之间的全部差异，其中既包括元素参考能不同带来的系统偏差，也包括目标数据真正要求模型学习的结构间能量变化。在只有能量标签的情况下，训练只要求模型把各个结构的能量预测准确，却不会告诉模型哪些误差只来自元素参考能差异。因此，模型会在降低能量误差的同时改变不同结构之间的相对能量和原子力。提前处理元素参考能后，这部分差异不再需要由网络本身承担，模型主要学习剩余的能量差异。改变学习率以及允许元素参考能继续随训练调整后，两种训练方式之间的差距仍然存在，说明这一结果并不是由某一个训练参数造成的。不同设置下，direct 与 aligned 之间的主要差距仍然存在，说明这一结果并不依赖某一特定学习率，也不是由元素参考能在训练中是否固定造成的。

真实形成能数据本身不能直接说明这些额外变化是否由元素参考能差异造成，因为目标数据可能确实要求不同结构之间的相对能量发生变化。为此，我们进一步构造了一个只改变元素参考能的受控任务。在这个任务中，目标能量只比预训练模型多出一个与元素组成有关、但与原子坐标无关的能量项，因此不同结构之间的相对能量和原子力理论上都不需要改变。即使如此，direct 仍然产生了明显的相对能量和原子力变化，而提前处理元素参考能后，这些变化接近正常训练带来的背景水平。进一步分析表明，direct 实际上已经基本学到了目标中的元素参考能变化，因此这些额外变化不能简单解释为参考能没有拟合好。这个实验说明，元素参考能差异本身就足以使仅用能量的微调产生目标并不需要的相对能量和原子力变化。

需要注意的是，在真实任务中，微调后变化越小并不一定越好。新的计算方法或新的能量数据本身可能要求不同结构之间的相对能量发生改变，因此预训练模型也不是始终应该保持不变。本文更关注的是，模型实际产生的变化是否与目标真正要求的变化一致。SevenNet 的结果表明，direct 产生了大量目标要求之外的变化，而提前处理元素参考能后这种偏离明显减小。固定的 Materials Project 弛豫轨迹也给出了相同趋势：aligned 对原有结构间能量关系的改变明显小于 direct，但并不意味着微调后的模型一定优于未经微调的预训练模型。
在真实任务中，微调后变化更小并不一定代表结果更好。不同的的计算方法或能量数据本身可能要求结构间相对能量发生改变，因此评价微调结果时，还需要判断这些变化是否符合目标数据的要求。SevenNet 中，direct 产生了较多目标要求之外的相对能量变化，而提前处理元素参考能后，这部分变化明显减少。对固定 Materials Project 弛豫轨迹的分析也得到类似结果：aligned 对原有结构间能量关系的改变小于 direct。

已有研究已经指出，元素参考能会影响预训练原子势在不同计算方法之间迁移时的能量误差、训练稳定性和最终预测结果\cite{huang_transfer}。本文关注的是，如果这部分差异没有提前处理，而是直接让模型在仅用能量标签的微调中拟合，观察模型还会同时改变什么。通过比较目标数据要求的相对能量变化与模型实际产生的变化，我们可以判断微调是否进行了额外调整；受控实验则进一步说明，单独的元素参考能差异已经足以引起这种现象。SevenNet、EquiformerV3 以及 GGA/GGA+$U$  到 r$^2$SCAN 的任务都得到相同趋势，这一结果并不只出现在某一个模型或某一套能量数据中。

这些结果对仅用能量标签的适配实际使用也有直接意义。这种微调有着充足的使用背景，如果目标任务主要关注相同化学组成下不同结构之间的能量排序，且预训练模型本身已经表现较好时，直接微调并不一定是最合适的选择；先检查并处理元素参考能差异，可以减少不必要的模型变化。如果目标任务还要求更准确的跨组成形成能，继续微调仍然有价值。本文中，只加入元素参考能修正后形成能 MAE 仍为 \(109.7~\mathrm{meV/atom}\)，而进一步微调可以降至 \(21.5~\mathrm{meV/atom}\)。提前处理元素参考能的目的不是阻止模型继续学习，而是避免把本来可以单独处理的能量差异也交给模型去拟合。评价微调效果时，也不应只看最终形成能 MAE，还应同时检查不同结构之间的相对能量是否发生了目标之外的改变；如果有可靠的原子力数据，还可以进一步检查微调前后的原子力变化。

本文仍有几个需要说明的限制。真实形成能数据没有对应的目标原子力，因此关于原子力变化的直接证据主要来自受控实验。部分评价数据参与过前期开发，并不是完全独立的盲测数据，这可能影响对模型最终泛化能力的判断。固定的 Materials Project 弛豫轨迹也可能与 SevenNet 的预训练数据存在重叠，因此这部分结果主要用于比较 direct 与 aligned 对已有结构间能量关系的影响。

后续可以在完全独立的数据上进一步验证本文结果，并尽可能排除与预训练数据的结构重叠。如果能够使用同时具有能量和原子力标签的数据，还可以直接检查真实适配任务中原子力是否也会受到类似影响，并比较提前处理元素参考能与加入力监督之间的关系。这样可以把目前主要依赖相对能量和受控实验得到的结论扩展到真实的能量—力适配任务。

另一个值得继续研究的问题是，这种额外变化在什么条件下会变得更明显。目标数据量、元素参考能差异大小、可训练参数范围以及模型结构都可能影响微调结果。本文已经在 SevenNet、EquiformerV3 和不同目标能量数据中观察到相同趋势，后续可以进一步研究这些因素如何影响模型对参考能差异和结构相关变化的分配，并据此形成更稳定的仅用能量标签的微调方法。

总体来看，预训练模型与新能量数据之间的差异并不都需要通过改变模型本身来学习。元素参考能的变化可以改变绝对形成能，却不要求化学组成相同的不同结构之间的相对能量和原子力随之改变。如果这部分差异直接进入仅用能量标签的微调，模型可能在降低形成能误差的同时产生额外的相对能量和原子力变化。提前处理元素参考能差异能够明显减少这种现象。因而，判断一次微调是否有效，不仅要看最终能量误差降低了多少，还需要判断模型为了得到这一精度究竟改变了什么，以及这些变化是否是目标数据真正需要的。

\section{方法}

\textbf{数据、模型与训练设置。}

主要实验使用 Matbench 的 \nolinkurl{matbench_mp_e_form} 数据集
\cite{dunn2020benchmarking}。该数据集整理自 Materials Project，提供晶体结构及对应的单位原子形成能。我们以预训练 SevenNet 的 \nolinkurl{7net-0} 模型为起点，对这套形成能数据进行适配\cite{park2024scalable}。EquiformerV3（EqV3）实验使用 \nolinkurl{omat24-mptrj-salex_gradient.pt} 作为初始模型\cite{liao2026equiformerv3}。两种模型均使用其实际对应的物理总能量输出进行训练，不根据预训练模型的预测重新定义目标标签。

SevenNet 主要实验采用固定的数据划分 \nolinkurl{medium_same_split_269}，其中训练集、验证集和原始测试集分别包含 8192、1024 和 2048 个结构。训练集用于模型训练和元素参考能拟合，验证集用于选择模型。原始测试集主要用于后面的受控实验和原子力评价。真实形成能任务的绝对能量 MAE 在另一套包含 2048 个结构的固定评价集上计算，该集合与原始测试集没有相同的结构 ID。结构间相对能量的评价使用另一套固定数据，其中包含 401 个结构、190 组元素种类和数量完全相同的结构，以及 238 个固定结构对。两套评价数据均参加过前期开发，因此本文将其作为固定评价集使用，而不视为完全独立的盲测集。为保证不同训练方式在相同的数据划分下进行比较，本文使用上述固定划分，而不采用 Matbench 官方五折交叉验证。主要结果使用三个固定随机种子：20260528、20260529 和 20260530。

SevenNet 训练时开放最后两个 interaction blocks，以及
\nolinkurl{reduce_input_to_hidden}、
\nolinkurl{reduce_hidden_to_energy} 和
\nolinkurl{rescale_atomic_energy}，
共 301300 个可训练参数。优化器采用 AdamW\cite{loshchilov2017decoupled}，学习率为 \(3\times10^{-4}\)，weight decay 为 \(10^{-5}\)，训练和验证的 batch size 分别为 16 和 32。损失函数为单位原子能量上的 Smooth L1，\(\beta=0.05~\mathrm{eV/atom}\)。主要实验训练 120 轮，不使用 early stopping、学习率调度、梯度裁剪或自动混合精度。

为了检查结果是否依赖具体训练设置，我们另外使用 \(10^{-4}\)、\(3\times10^{-4}\) 和 \(10^{-3}\) 三个学习率进行对照。在另一组实验中，训练前拟合得到的元素参考能不再固定，而是允许其随模型一起训练。除这些变化外，其余训练设置均与主要实验相同。

EqV3 使用 LoRA 对主干网络进行适配\cite{hu2021lora}，同时训练完整的顶层能量模块 \nolinkurl{energy_block}。LoRA 加在 \nolinkurl{blocks.*} 下的 \nolinkurl{torch.nn.Linear} 和 \nolinkurl{SO3Linear} 权重上，rank 为 8，\(\alpha=16\)，不使用 dropout。LoRA 和能量模块共包含 860417 个可训练参数。学习率为 \(2\times10^{-4}\)，AdamW 的 weight decay 为 \(10^{-4}\)，batch size 为 8，梯度范数裁剪阈值为 2.0。训练使用 ReduceLROnPlateau 调整学习率，factor 为 0.5，patience 为 20，并采用 bfloat16 autocast。EqV3 共训练 40 轮，以验证集单位原子 MAE 最低的最早轮次作为最终模型。

若 EqV3 网络输出记为 \(q\)，其物理总能量写为

\[
E_{\mathrm{phys}}
=
\sigma q+m+\sum_Z n_Z\epsilon_Z+N_Xh(z),
\]

其中 \(\sigma=0.6274358034\)、\(m=0\)，\(h(z)\) 为下文定义的、只与元素组成有关的参考能项。direct 中 \(h(z)=0\)，提前处理元素参考能的训练方式则加入对应的 \(h(z)\)。训练目标为 \(N_Xy\)，并在总能量上使用 Smooth L1 损失，\(\beta=0.05~\mathrm{eV}\)。

GGA/GGA+$U$  到 r$^2$SCAN的实验沿用 SevenNet 的可训练模块和主要训练设置。目标数据包含 10000 个 Materials Project r$^2$SCAN 记录，训练集、验证集和测试集分别包含 6915、1431 和 1654 个结构。元素参考能只使用训练集拟合。该任务训练 60 轮，其余设置与 SevenNet 主要实验保持一致。

为了检查微调后已有结构之间的能量关系，我们还使用了 Materials Project 公开的 GGA/GGA+$U$ 弛豫轨迹。所有轨迹在进行模型预测之前完成筛选并固定，不再执行新的结构弛豫。最终数据包含 50 个计算任务、171 个去重结构和 434 对同一任务中的结构，其中 37 个任务采用 GGA，13 个采用 GGA+$U$。该实验只比较同一弛豫任务中不同结构之间的相对能量。由于无法确定这些结构是否曾出现在 SevenNet 的原始预训练数据中，因此该实验主要用于比较不同训练方式对这些固定结构相对能量的影响。

实验环境为 Python 3.11.15、PyTorch 2.7.1+cu118、SevenNet 0.13.0、ASE 3.28.0、scikit-learn 1.8.0 和 NumPy 2.4.6，计算节点使用 4 张 NVIDIA GeForce RTX 4090。EqV3 使用固定 checkpoint 和训练时保存的源码版本。


\textbf{元素参考能与两种训练方式。}

记结构 \(X\) 的总能量为 \(E_\theta(X)\)，原子数为 \(N_X\)，对应的单位原子能量为

\[
e_\theta(X)=\frac{E_\theta(X)}{N_X}.
\]

对于元素 \(Z\)，其在结构中的原子分数记为

\[
x_Z=\frac{n_Z}{N_X}.
\]

我们用一个只依赖元素组成的线性项表示训练数据与预训练模型之间的元素参考能差异：

\[
h(z)=b_0+\sum_Z x_Zb_Z.
\]

实际计算使用 118 个元素原子分数和一个截距。SevenNet 中，这一项只使用 8192 个训练结构拟合。具体做法是根据预训练模型与目标形成能之间的残差进行 Ridge 回归，正则化系数为 \(10^{-6}\)。验证集和最终评价集都不参与元素参考能的拟合。

direct 直接使用模型原来的单位原子能量预测目标形成能：

\[
\hat e_{\mathrm{direct}}(X)
=
e_\theta(X).
\]

另一种训练方式先将训练集拟合得到的元素参考能加入模型输出：

\[
\hat e_{\mathrm{aligned}}(X)
=
e_\theta(X)+h(z).
\]

两种训练方式真正的区别只是训练前是否先处理元素参考能差异。由于 \(h(z)\) 只与元素组成有关，而不依赖原子坐标，

\[
\nabla_{\mathbf R}[N_Xh(z)]=0,
\]

所以这一项本身不会产生原子力，也不会改变化学组成相同结构之间的能量差。

Frozen 表示完全不继续训练的预训练 SevenNet；Frozen+$h$ 仍然保持 SevenNet 不变，只在其预测结果上加入前面拟合出的元素参考能。这样可以直接判断，仅处理元素参考能差异本身能够降低多少形成能误差。

为了保证 direct 和 aligned 的比较公平，训练结束后，两种方法都对训练后模型预测中仍然存在的、只与元素组成有关的剩余误差进行一次相同的线性修正：

\[
\hat e^{\mathrm{cal}}
=
\hat e^{\mathrm{model}}
+
x^{\mathsf T}w+b.
\]

其中 \(w\) 和 \(b\) 只使用训练集残差拟合，Ridge 系数同样为 \(10^{-6}\)。这一修正不参与模型训练，也不会改变化学组成相同结构之间的能量差。


\textbf{只改变元素参考能的受控实验。}

真实形成能数据中，不同结构之间的相对能量本身可能需要发生变化，因此真实任务无法单独分离元素参考能差异对微调结果的影响。为此，我们另外构造了一个只改变元素参考能的任务。

以冻结 SevenNet 的总能量 \(E_0(X)\) 为起点，定义新的目标能量

\[
E_{\mathrm{syn}}(X)
=
E_0(X)+N_Xh(z).
\]

这里的 \(h(z)\) 与 SevenNet 真实任务中使用的元素参考能函数完全相同。对于元素种类和数量完全相同的两个结构 \(a\) 和 \(b\)，

\[
E_{\mathrm{syn}}(a)-E_{\mathrm{syn}}(b)
=
E_0(a)-E_0(b).
\]

同时，由于 \(h(z)\) 不依赖原子坐标，

\[
-\nabla_{\mathbf R}E_{\mathrm{syn}}(X)
=
-\nabla_{\mathbf R}E_0(X).
\]

因此，这个受控目标只改变绝对能量，目标中的结构间相对能量和原子力与 Frozen 完全相同。

我们设置三种训练方式。zero\_direct 直接重新拟合预训练模型自身的能量，用于测量正常训练带来的背景变化；reference\_direct 直接拟合加入元素参考能后的目标；reference\_aligned 使用相同的目标，但在训练前先加入这部分元素参考能。三种训练使用相同的数据划分、优化器、损失函数、120 轮训练预算和三个随机种子。


\textbf{结构间相对能量的评价。}

为了排除元素参考能的共同影响，我们只在元素种类和数量完全相同、但原子排列不同的结构之间比较相对能量。仅具有相同约化化学式并不足以进入同一组，只有完整元素计数向量完全相同的结构才进行比较。

对于一对结构 \(a\) 和 \(b\)，其结构间能量差预测误差定义为

\[
\left|
(\hat e_a-\hat e_b)
-
(y_a-y_b)
\right|.
\]

在所有固定结构对上取平均，得到正文使用的结构间能量差 MAE：

\[
\Delta E\text{-MAE}
=
\frac{1}{|\mathcal P|}
\sum_{(a,b)\in\mathcal P}
\left|
(\hat e_a-\hat e_b)
-
(y_a-y_b)
\right|.
\]

我们还统计每一组中预测能量最低的结构是否与真实最低能结构一致，并将全部组中的命中比例作为最低能结构命中率。若存在数值并列的真实最低能结构，只要预测结果命中其中任意一个均记为正确。

为了进一步判断微调产生的变化是否是目标数据真正要求的，记预训练模型的单位原子能量为 \(e_0\)，目标形成能为 \(y\)，微调后的预测为 \(e_s\)。对于一对元素组成相同的结构 \(a\) 和 \(b\)，目标数据要求的能量差变化定义为

\[
t_{ab}
=
(y_a-y_b)
-
(e_{0,a}-e_{0,b}),
\]

模型经过微调后实际产生的变化定义为

\[
c_{s,ab}
=
(e_{s,a}-e_{s,b})
-
(e_{0,a}-e_{0,b}).
\]

若

\[
c_{s,ab}=t_{ab},
\]

说明模型对这一对结构产生的变化与目标完全一致。

将所有结构对上的 \(t_{ab}\) 和 \(c_{s,ab}\) 分别组成向量 \(t\) 和 \(c_s\)。模型产生的总体变化与目标要求之间的比例写为

\[
\frac{\|c_s\|}{\|t\|}.
\]

为了进一步区分模型沿目标要求产生的变化和额外变化，我们计算

\[
\alpha_s
=
\frac{c_s^{\mathsf T}t}{t^{\mathsf T}t}.
\]

其中 \(\alpha_s t\) 表示模型变化中沿目标要求方向的部分，而

\[
\frac{\|c_s-\alpha_s t\|}{\|t\|}
\]

表示不能由目标要求的变化解释的部分。若模型实际产生的变化与目标完全一致，则总体变化比和 \(\alpha_s\) 均为 1，而目标要求之外的部分为 0。

我们还直接计算模型实际变化与目标要求之间的距离：

\[
\frac{\|c_s-t\|}{\|t\|}.
\]

如果模型完全保持预训练结果不变，则 \(c_s=0\)，该值等于 1。因此，小于 1 表示微调后的变化比完全保持预训练模型更接近目标，大于 1 则表示更加偏离目标。


\textbf{受控实验中的参考能学习与原子力评价。}

在只改变元素参考能的受控实验中，我们还检查 direct 是否真正学到了人为加入的元素参考能。记训练后模型相对于 Frozen 的单位原子能量变化为向量 \(u\)，人为加入的元素参考能变化为 \(h\)，计算

\[
\beta
=
\frac{u^{\mathsf T}h}{h^{\mathsf T}h}.
\]

当 \(\beta\) 接近 1 时，说明模型已经基本重现人为加入的参考能变化。扣除这一部分以后，剩余变化写为

\[
u-\beta h,
\]

并用

\[
\frac{\|u-\beta h\|}{\|h\|}
\]

比较 direct 和 aligned 中额外能量变化的大小。

原子力由模型总能量对原子坐标求导得到：

\[
\mathbf F
=
-\nabla_{\mathbf R}E.
\]

在受控任务中，元素参考能本身不会产生原子力，因此 Frozen 的原子力就是这一任务的目标原子力。我们用训练后模型与 Frozen 之间的力分量 RMSE 衡量微调带来的原子力变化。

最初的原子力评价使用原始测试集中的 12 个固定结构，并进一步扩展到 100 个结构检查结果是否由少数样本造成。结构的选择在模型预测之前完成，不使用真实标签或模型结果。我们还使用中心差分检查能量与原子力计算的一致性。


\textbf{统计分析}

除特别说明外，正文结果均报告三个固定随机种子的均值和样本标准差。置信区间采用 10000 次 bootstrap 重采样计算。

对于元素组成相同的不同结构，以完整结构组作为重采样单位；原子力评价以完整结构作为重采样单位；Materials Project 弛豫轨迹则以完整计算任务作为重采样单位。不同训练方式之间的比较使用相同的重采样样本，以保持配对关系。

95\% 置信区间取 bootstrap 分布的 2.5\% 和 97.5\% 分位数。区间不跨 0 表示在当前数据中两种训练方式的差异方向明确；区间跨 0 则表示当前数据不足以确定差异方向。


\subsection{数据与代码可用性}

本文实验使用的数据划分、预测结果、固定结构对、bootstrap 结果、原子力评价、有限差分检查、训练过程记录、学习率对照、元素参考能对照以及 Materials Project 弛豫轨迹清单均已保存。相关数据和代码尚未公开发布。

\clearpage
\bibliographystyle{unsrt}
\bibliography{references}

\end{document}